\documentclass[12pt]{article}
\usepackage[utf8]{inputenc}
\usepackage[T1]{fontenc}
\usepackage{amsmath, amssymb}
\usepackage{hyperref}

\title{Introdução à Análise de Fourier}
\author{Seu Nome}
\date{\today}

\begin{document}
\maketitle

\begin{abstract}
Este artigo apresenta os conceitos fundamentais da análise de Fourier,
incluindo séries e transformadas, com aplicações em sinais e sistemas.
\end{abstract}

\section{Séries de Fourier}

Uma função periódica $f(x)$ com período $2\pi$ pode ser escrita como:

\[
f(x) = \frac{a_0}{2} + \sum_{n=1}^{\infty} \left( a_n \cos(nx) + b_n \sin(nx) \right)
\]

onde os coeficientes são dados por:

\[
a_n = \frac{1}{\pi} \int_{-\pi}^{\pi} f(x)\cos(nx)\,dx, \qquad
b_n = \frac{1}{\pi} \int_{-\pi}^{\pi} f(x)\sin(nx)\,dx
\]

\subsection{Exemplo: Onda Quadrada}

Para a onda quadrada $f(x) = \text{sgn}(\sin x)$, os coeficientes valem:

\[
b_n = \begin{cases} \dfrac{4}{n\pi} & \text{se } n \text{ é ímpar} \\ 0 & \text{se } n \text{ é par} \end{cases}
\]

\section{Transformada de Fourier}

Para funções não periódicas, usamos a transformada:

\[
\hat{f}(\xi) = \int_{-\infty}^{\infty} f(x)\, e^{-2\pi i x \xi}\, dx
\]

E a transformada inversa:

\[
f(x) = \int_{-\infty}^{\infty} \hat{f}(\xi)\, e^{2\pi i x \xi}\, d\xi
\]

\subsection{Propriedades Importantes}

\begin{itemize}
  \item \textbf{Linearidade:} $\widehat{af + bg} = a\hat{f} + b\hat{g}$
  \item \textbf{Deslocamento:} $\widehat{f(x-a)}(\xi) = e^{-2\pi i a\xi}\hat{f}(\xi)$
  \item \textbf{Convolução:} $\widehat{f * g} = \hat{f} \cdot \hat{g}$
  \item \textbf{Parseval:} $\displaystyle\int |f|^2\,dx = \int |\hat{f}|^2\,d\xi$
\end{itemize}

\section{Conclusão}

A análise de Fourier é uma ferramenta central em matemática aplicada,
processamento de sinais, física e engenharia.

\end{document}
